(1) 刚体绕固定 \(z\) 轴以角速度 \(\Omega\) 旋转，磁偶极矩大小 \(m\)，与自转轴夹角 \(\alpha\) 固定。写出 \(\boldsymbol m(t)\) 并由远区磁偶极辐射公式 \(P=\mu_0|\ddot{\boldsymbol m}|^2/(6\pi c^3)\) 推导周期平均辐射功率，检查 \(\alpha=0\) 极限。
(2) 刚体绕轴转动惯量 \(I\) 恒定，假设全部辐射能量由自转动能供给，忽略磁矩和倾角演化。推导制动方程并精确求 \(\Omega(t)\)，给出初始慢变条件。
(3) 定义制动指数 \(n=\Omega\ddot\Omega/\dot\Omega^2\) 与自转周期 \(P_{\rm rot}=2\pi/\Omega\)。求 \(n\) 及 \(P_{\rm rot}\dot P_{\rm rot}\) 的常数表达式，讨论与观测偏离的可能物理来源。
(4) 求从 \(\Omega_0\) 降到 \(\Omega_1\) 的辐射总能量和时间，核对角动量损失的量纲与能量守恒，说明理想真空偶极制动对真实磁层的局限。
